% 4. Arguing Against Myself
\section{Arguing Against Myself}
\label{sec:against}

So far I have made the case as well as I can. Now I want to attack it. Below
are the eight strongest objections I know of. For each, I try to state it as
its best advocate would, and then say whether I think it can be answered or
whether I have to concede. Some of these I concede entirely.

\subsection{Survivorship bias}

\emph{The objection.} Every case in this paper was chosen because it worked.
For each bottle light there are thousands of dark rooms; for each M-Pesa there
are mobile-money launches that never took off. If you only look at the
winners, any condition they share will look like a cause of winning, even if
it is equally common among the losers. Scarcity is extremely common. Seeing it
behind every frugal success tells us nothing unless we also know how often
scarcity produced nothing.

\emph{My answer.} I concede most of this. I do not have a census of failed
frugal innovations, and I do not think one exists. \todo{NeoNurture} The honest
response is to stop claiming that scarcity \emph{causes} innovation and to
claim something weaker: that when innovation does happen under scarcity, it
tends to take a recognisable form, cheaper and narrower and built around the
missing system. That claim is about the \emph{kind} of solution, not its
frequency, and survivorship bias does not touch it in the same way. The
frequency question stays open, and I return to it in \cref{sec:conclusion}.

\subsection{Hidden abundance}

\emph{The objection.} The solar kit is not a product of scarcity. It is a
product of decades of semiconductor research, a global photovoltaic industry
and the lithium-ion battery, all developed in rich countries for rich-country
markets and made cheap by enormous investment. What happened in Kenya was the
last step, distribution, and calling that ``scarcity innovation'' takes
credit for an iceberg because you can see its tip.

\emph{My answer.} This is true, and \cref{fig:hidden-abundance} tries to show
how true. \todo{HiddenAbundance} But I think it changes the thesis rather than
refuting it. The question was never whether scarcity can produce solar cells.
It cannot. It was whether the absence of a grid produces a different
\emph{use} of solar cells than its presence does, and on that the evidence is
reasonably clear: the off-grid kit with a pay-as-you-go plan is not a product
anyone designed for a house with a meter. Abundance made the parts; scarcity
chose the assembly. That is the synthesis I develop in
\cref{sec:synthesis}.

\begin{figure}[htbp]
\centering
% Figure: the hidden abundance behind a solar kit. Shares are China's share of
% global manufacturing capacity in 2021 (data/manual/literature_figures.csv).
\begin{tikzpicture}[
  node distance=4mm and 10mm,
  up/.style={concept, text width=0.36\linewidth, align=flush left, font=\sffamily\footnotesize},
  stat/.style={font=\sffamily\footnotesize\bfseries, text=indigo},
]
  \node[up] (pv) {\textbf{Solar PV manufacturing}\\[2pt]
    China's share of capacity: polysilicon \pvPolysiliconShare{},
    wafers \pvWaferShare{}, cells \pvCellShare{}, modules \pvModuleShare{}};
  \node[up, below=of pv] (batt) {\textbf{Battery manufacturing}\\[2pt]
    \batteryCapacityShare{} of global battery production capacity in China,
    built for electric vehicles};
  \node[up, below=of batt] (chips) {\textbf{LEDs and control chips}\\[2pt]
    Semiconductor fabs and decades of lighting research};
  \node[up, below=of chips] (net) {\textbf{Mobile networks and mobile money}\\[2pt]
    Towers, handsets, SIMs and payment rails, for pay-as-you-go};

  \node[concept accent, text width=0.2\linewidth, right=14mm of $(batt.east)!0.5!(chips.east)$] (kit)
    {\textbf{Solar home kit}\\[2pt] panel, battery, lights, phone charger};
  \node[concept counter, text width=0.17\linewidth, right=8mm of kit] (home)
    {\textbf{Off-grid household}\\[2pt] pays a little each week};

  \foreach \n in {pv, batt, chips, net} \draw[flow] (\n.east) -- (kit.west);
  \draw[flow, draw=clay] (kit) -- (home);

  \node[note, below=3mm of $(kit.south)!0.5!(home.south)$, text width=0.4\linewidth] {Assembled, sold and financed under scarcity};
  \node[note, anchor=north west, text width=0.36\linewidth, align=left] at ($(net.south west)+(0,-2mm)$)
    {Built by abundance. Module costs fell \moduleCostFall{} between 2009 and 2023.};
\end{tikzpicture}

\caption{\takeaway{A cheap solar lantern rests on expensive, concentrated
abundance.} The kit sold to an off-grid household is the last link in a chain
of capital-intensive industries, most of them concentrated in a few countries
and none of them built for off-grid markets.
\figsource{\todo{HiddenSource}}}
\label{fig:hidden-abundance}
\end{figure}

\subsection{The limits of the fixes}

\emph{The objection.} Even the successes are partial. The bottle light gives
no light at night and needs a thin roof and a sunny day. Most off-grid solar
users have small lights that fall below the lowest tier of electricity access.
Mobile money moves money but does not, on its own, give credit or insurance on
fair terms. A frugal solution that does a tenth of what the original does is
not a different kind of more; it is less.

\emph{My answer.} Partly conceded. \todo{Limits} I have tried to report these
limits in each case study, and I think the right summary is that frugal
solutions are \emph{narrower} than the systems they stand in for. Sometimes
narrowness is exactly what makes them affordable, and for the household that
gets light for the first time, a tenth of a grid is a great deal more than
nothing. But nobody should mistake a narrower solution for the full one, and a
policy that treats the solar lantern as ``job done'' would be making a serious
mistake.

\subsection{Romanticising hardship}

\emph{The objection.} Stories of ingenious poor people are pleasant to read in
rich countries. They flatter the reader's sense that poverty is a kind of
character-building, and they can quietly excuse the absence of investment: if
villagers can light their homes with bottles, why build them a grid?
\textcite{birtchnell2011} makes a version of this argument about jugaad, which
he reads as a symptom of systemic failure rather than a model to admire.

\emph{My answer.} This is the objection I take most seriously, because the
risk is real and the argument in this paper could be misused in exactly this
way. My response is not to deny it but to state the position plainly. Nothing
here suggests that scarcity is good for people. The psychology in
\cref{sec:psychology} says the opposite: deprivation narrows attention and
damages decisions. The conditions in \cref{sec:synthesis} under which scarcity
sometimes produces something useful, tools, slack, room to experiment and
paths to scale, are all things that investment provides. If this paper has a
policy implication, it is to invest more, not less, in the places where
frugal ideas come from.

\subsection{Attribution and outside capital}

\emph{The objection.} M-Pesa was built by Safaricom, a company part-owned by
Vodafone, using a grant from the UK government. Aravind's lenses rely on
technology from rich-country medicine. Liter of Light is a foundation. Calling
these ``local'' innovations born of scarcity hides who paid for them and who
profits from them.

\emph{My answer.} Partly conceded, and worth separating carefully. In the
M-Pesa case, the local part was the need and the use: Kenyans repurposed a loan
repayment pilot into a way of sending money home \parencite{hughes2007}. The
outside part was the capital and the platform: the Department for
International Development contributed \dfidGrant{} and Vodafone matched it
with \vodafoneMatch{} \parencite{phillips2013dfid}. Both were necessary. The
mechanism in \cref{fig:mechanism} survives, but its third condition, a path to
scale, turns out to depend heavily on outside resources. That is a large part
of why the synthesis treats scarcity and abundance as two modes of one system
rather than as alternatives.

\subsection{Measurement}

\emph{The objection.} ``Great innovation'' is undefined. By what measure is a
bottle light more impressive than a power station? If you measure impact per
dollar, frugal solutions win almost by definition; if you measure total impact,
the grid wins easily. The choice of metric decides the argument.

\emph{My answer.} This is right, and the honest thing is to show both rankings
(\cref{tab:measurement}). Ranked by cost per unit, the bottle light and the
Jaipur Foot come first. Ranked by people reached, mobile money and off-grid
solar come first, and the bottle light near the bottom. Neither ranking is
wrong. The thesis is that scarcity produces solutions that do well on the
first measure, and some of which, under the right conditions, also do well on
the second. It is not that frugal solutions are better on every measure.

\begin{table}[htbp]
\centering
\caption{\takeaway{Cost per unit and total reach rank the cases differently.}
The cheapest solutions per unit are not the ones that reach the most people.
Figures are not comparable across rows in year or currency; they show orders of
magnitude, not a league table.}
\label{tab:measurement}
\small
\begin{tabularx}{\linewidth}{@{}l X X@{}}
\toprule
Case & Cost per unit & Reach \\
\midrule
Bottle light & Near zero: a bottle, water and bleach & \literHomesPhilippines{} homes in the Philippines (2013) \parencite{zobel2013bbc} \\
Jaipur Foot & About \jaipurCostOrig{} a foot (1978) \parencite{sethi1978} & \todo{JaipurReach} \\
Embrace warmer & \embraceCostLow{}--\embraceCostHigh{} (2011 estimate) \parencite{routson2011embrace} & No published count \\
Aravind cataract surgery & \aravindCostPerOp{} an operation \parencite{le2016aravind} & More than \aravindSurgeries{} surgeries and laser procedures \parencite{ravilla2014} \\
Off-grid solar, SSA & Varies by kit; mostly paid in instalments & \sspSolarUsers{} users (\sspSolarYear{}) \parencite{irena2025offgrid} \\
Mobile money & Free to open; fees per transaction & \kenMobileMoney{} of Kenyan adults; \gsmaRegisteredAccounts{} accounts worldwide \parencite{worldbank_findex,gsma2026sotir} \\
\bottomrule
\end{tabularx}
\end{table}

\subsection{The inverted U}

\emph{The objection.} By the constraints literature's own account, severe
scarcity hurts creativity \parencite{acar2019}, and the places in this paper
are not facing mild constraints. If the inverted U is right, the thesis
applies to the wrong people: it describes what happens to moderately
constrained engineers, not to households without electricity.

\emph{My answer.} I accept the inverted U and think it sharpens the thesis.
The key is to separate the constraint on the \emph{system} from the constraint
on the \emph{person}. The bottle light, M-Pesa and the solar kit all emerged
where a system was missing entirely but the people designing and adopting them
still had some slack: a mechanic with a roof and tools, a mobile operator with
capital, households with enough income to pay a little each week. Where
scarcity reaches the person, so that there is no slack left at all, I expect
nothing, and the evidence agrees \parencite{mani2013}.

\subsection{Causation}

\emph{The objection.} Suppose mobile money and off-grid solar did spread fastest
where banks and grids were weakest. That correlation could have many causes
besides scarcity: liberal telecoms regulation, a single dominant operator, high
phone ownership, cheap Chinese components, donor money. Scarcity might be a
bystander.

\emph{My answer.} I cannot rule this out, and I have not tried to. Nothing in
this paper is a causal estimate of the effect of scarcity on innovation. The
cases are illustrations, and the cross-country data in
\cref{fig:electricity-access,fig:findex-accounts} show what happened, not why.
\todo{Causation} The most I can say is that the mechanism is plausible, that it is
consistent with the order of events in each case, and that it predicts
something the alternatives do not: the \emph{form} the solutions took. A
liberal regulator explains why mobile money was allowed. It does not explain
why it was designed around corner shops rather than branches. That is what the
missing default explains. But I hold this loosely, and \cref{sec:conclusion}
says what evidence would change my mind.
